\section*{Chapter \ref{ch:b2:magmatic-differentiation} --- Magmatic Differentiation}

\begin{solution}{exo:b2:magmatic-differentiation:1}
(a) The liquidus of An$_{50}$ is \qty{1450}{\celsius}. The first crystal lies on the
solidus at that temperature, between An$_{66.7}$ (1372) and An$_{83.3}$ (1465):
$66.7 + 16.6\times(1450 - 1372)/(1465 - 1372) = 66.7 + 13.9 = \qty{80.6}{\percent}$
anorthite. (b) At \qty{1394}{\celsius} the liquid is on the liquidus,
An$_{33.3}$; the crystals on the solidus, $66.7 + 16.6\times22/93 =
\qty{70.6}{\percent}$. By the lever rule, the crystallised fraction is $(50 -
33.3)/(70.6 - 33.3) = 16.7/37.3 = 0.448$. (c) The last liquid disappears when the
crystals reach An$_{50}$, at the solidus temperature \qty{1287}{\celsius}; the
liquidus there lies between An$_{11.1}$ (1265) and An$_{20}$ (1334): $11.1 +
8.9\times22/69 = \qty{13.9}{\percent}$ anorthite.
\end{solution}

\begin{solution}{exo:b2:magmatic-differentiation:2}
$\bar D_{\ce{Sr}} = 0.15\times0.000063 + 0.35\times0.157 + 0.50\times2.7 = 0.00001 +
0.0550 + 1.35 = 1.40$; $\bar D_{\ce{Zr}} = 0.0001 + 0.0683 + 0.0005 = 0.069$;
$\bar D_{\ce{Rb}} = 0.00001 + 0.0012 + 0.0125 = 0.014$. Only strontium, held by the
plagioclase, is compatible ($\bar D > 1$); zirconium and rubidium are incompatible,
rubidium the more so.
\end{solution}

\begin{solution}{exo:b2:magmatic-differentiation:3}
(a) $C_L/C_0 = 0.5^{D-1}$: Sr $0.5^{0.40} = 0.758$ (it falls); Zr $0.5^{-0.931} =
1.907$; Rb $0.5^{-0.986} = 1.981$, close to the limit 2. (b) Zr/Rb changes by
$0.5^{0.069 - 0.014} = 0.5^{0.055} = 0.963$, less than \qty{4}{\percent} while each
element nearly doubled: as the corollary says, the ratio of two \omterm{def:b2:elements-in-the-earth:partition}{incompatible
elements} keeps the value of the parent \omterm{def:b2:magma-emplacement:magma}{magma}, and so tells about the source, not
about crystallisation.
\end{solution}

\begin{solution}{exo:b2:magmatic-differentiation:4}
(a) \ce{SiO2} $+5.03$, \ce{MgO} $-1.95$, \ce{K2O} $+0.82$ (per cent), potash
multiplied by $1.67/0.85 = 1.96$. (b) For a perfectly \omterm{def:b2:elements-in-the-earth:partition}{incompatible element}
$C_L/C_0 = 1/F$, so $F = 0.85/1.67 = 0.509$: half of the \omterm{def:b2:magma-emplacement:magma}{magma} would have had to
crystallise. (c) Removing \qty{12.5}{g} per \qty{100}{g} leaves $F \approx 0.88$ and
would raise potash only by about $1/0.88 = 1.14$: crystallisation alone cannot
explain the rise. The potash came from outside, from the granitic xenoliths, rich in
potash (\qty{3.6}{\percent}), that the \omterm{def:b2:magma-emplacement:magma}{magma} assimilated.
\end{solution}

\begin{solution}{exo:b2:magmatic-differentiation:5}
(a) $F = 0.8$, $D = 5$. Fractional: $0.8^{4} = 0.410$; equilibrium:
$1/(5 + 0.8\times(1 - 5)) = 1/1.8 = 0.556$. (b) $F = 0.1$, $D = 0.01$. Fractional:
$0.1^{-0.99} = 9.77$; equilibrium: $1/(0.01 + 0.1\times0.99) = 1/0.109 = 9.17$.
(c) As $F \to 0$, the Rayleigh law tends to $0$ for $D > 1$ and to infinity for
$D < 1$, while the equilibrium law tends to $1/D$ (the last liquid in equilibrium
with crystals of the bulk composition). Only \omterm{def:b2:magmatic-differentiation:crystallisation}{fractional crystallisation} can strip a
liquid of a \omterm{def:b2:elements-in-the-earth:partition}{compatible element} completely.
\end{solution}

\begin{solution}{exo:b2:magmatic-differentiation:6}
(a) One formula unit holds 1.6 \ce{MgO}, 0.4 \ce{FeO} and one \ce{SiO2}: $1.6\times
40.30 = 64.48$, $0.4\times71.84 = 28.74$, 60.08, total \qty{153.30}{g}. Mass
percentages \ce{MgO} \qty{42.06}{\percent}, \ce{FeO} \qty{18.75}{\percent},
\ce{SiO2} \qty{39.19}{\percent}. (b) Lever rule on \ce{MgO}: $x = (10.0 - 8.0)/(42.06
- 8.0) = 2.0/34.06 = 0.0587$, about \qty{6}{\percent} of olivine. (c) Silica of the
liquid: $(49.0 - 0.0587\times39.19)/(1 - 0.0587) = 46.70/0.9413 =
\qty{49.61}{\percent}$. Removing olivine, poorer in silica than basalt, raises the
silica of the liquid, but only slowly.
\end{solution}

\begin{solution}{exo:b2:magmatic-differentiation:7}
(a) $F(0) = 0.768$; $F(1000) = 0.1091 - 0.5906 + 0.7678 = 0.286$; $F(2000) = 0.4364
- 1.1813 + 0.7678 = 0.023$. (b) The discriminant is $(5.9064\times10^{-4})^2 -
4\times1.091\times10^{-7}\times0.7678 = 3.4885\times10^{-7} - 3.3507\times10^{-7} =
1.379\times10^{-8}$, of square root $1.174\times10^{-4}$; the smaller root is
$H = (5.9064 - 1.174)\times10^{-4}/(2.182\times10^{-7}) = \qty{2169}{m}$, the
Sandwich Horizon. (c) $1 - 0.768 = 0.232$: nearly a quarter of the \omterm{def:b2:magma-emplacement:magma}{magma} had
crystallised before the base of the exposed Lower Zone formed, in the Hidden Zone
beneath it and in the Border Series along the walls and the roof.
\end{solution}

\begin{solution}{exo:b2:magmatic-differentiation:8}
The \omterm{def:b2:magmatic-differentiation:cumulate}{cumulate} holds $D\,C_L = D\,C_0F^{D-1}$, so between two levels
$\ln(\mathrm{Ni}_2/\mathrm{Ni}_1) = (D - 1)\ln(F_2/F_1)$. (a) $D - 1 = \ln(54/177)/
\ln(0.29/0.70) = (-1.187)/(-0.881) = 1.35$, so $D \approx 2.35$ (the regression over
the 66 \omterm{def:b2:magmatic-differentiation:cumulate}{cumulates} with $F \ge 0.23$ gives 2.38). (b) $D - 1 = \ln(6.2/54)/\ln(0.17/0.29) = (-2.164)/(-0.534)
= 4.05$, so $D \approx 5.05$. (c) The bulk $D$ changes when the crystallising
assemblage changes (a new mineral that takes up nickel, or one that leaves), and
mineral--melt coefficients themselves rise as the liquid cools and changes
composition, as measured for olivine; trapped liquid in the \omterm{def:b2:magmatic-differentiation:cumulate}{cumulates} blurs the
estimate.
\end{solution}

\begin{solution}{exo:b2:magmatic-differentiation:9}
(a) $100 + 25.4 - 2.9 - 9.6 = \qty{112.9}{g}$. (b) Grams of \ce{SiO2}: $55.53 +
25.4\times0.7248 - 2.9\times0.3919 - 9.6\times0.5054 = 55.53 + 18.41 - 1.14 - 4.85 =
\qty{67.95}{g}$, so $67.95/112.9 = \qty{60.19}{\percent}$ against 60.16. \ce{MgO}:
$4.66 + 0.20 - 1.22 = \qty{3.64}{g}$, \qty{3.22}{\percent} against 3.19. \ce{K2O}:
$0.88 + 0.92 = \qty{1.80}{g}$, \qty{1.60}{\percent} against 1.75. (c) Silica and magnesia, two of the three
controls, fit to within 0.03; potash, not used as a control, is low by 0.15, nearly a
tenth of its value: the assimilated material had more potash than the average
xenolith, or a potash-rich melt of it was taken preferentially.
\end{solution}

\begin{solution}{exo:b2:magmatic-differentiation:10}
(a) $z = (0.2 + 1.5 - 1)/(0.2 - 1) = 0.7/(-0.8) = -0.875$; $F^{-z} = 0.7^{0.875} =
0.7319$; $C^*/C_0 = 0.2\times0.5/(0.2 - 1 + 1.5) = 0.1/0.7 = 0.1429$, so $C/C_0 =
0.7319 + 0.1429\times(1 - 0.7319) = 0.7319 + 0.0383 = 0.770$. (b) Rayleigh:
$0.7^{0.5} = 0.837$. The AFC \omterm{def:b2:magma-emplacement:magma}{magma} has less strontium for two reasons: to lose
\qty{30}{\percent} of its mass while gaining $r$ times what it crystallises it must
crystallise $(1 - 0.7)/(1 - 0.2) = 0.375$ instead of 0.30, and the strontium it gains
comes from a rock half as rich. (c) $(0.7319\times0.7035 + 0.0383\times0.7200)/0.7702
= (0.5149 + 0.0276)/0.7702 = 0.70432$: the ratio has risen by 0.0008 while
crystallisation alone would have left it at 0.7035.
\end{solution}

\begin{solution}{exo:b2:magmatic-differentiation:11}
(a) $1.1\times(1100 - 200) + 250 = 990 + 250 = \qty{1240}{J}$ per gram. (b)
$25.4\times1240 = \qty{31496}{J}$, close to Wilcox's \qty{31350}{J}. (c)
$31496/370 = \qty{85.1}{g}$ of plagioclase. (d) The mass balance removed only $2.9 +
9.6 = \qty{12.5}{g}$ of crystals, whose latent heat (about \qty{4800}{J}) covers less
than a sixth of the need; $r = 25.4/12.5 = 2.03$, twice as much rock assimilated as
crystals formed, while the heat budget allows $r$ well below 1. The erupted \omterm{def:b2:magma-emplacement:magma}{magma}
cannot have supplied the heat itself: a larger mass of \omterm{def:b2:magma-emplacement:magma}{magma} must have crystallised
and cooled at depth, out of sight.
\end{solution}

\begin{solution}{exo:b2:magmatic-differentiation:12}
(a) Stage 1, $F$ from 1 to 0.8 with $D = 5$: factor $0.8^{4} = 0.4096$. Stage 2,
from 0.8 to 0.5 with $D = 1.5$: the law applies to the liquid remaining, so the
factor is $(0.5/0.8)^{0.5} = 0.7906$. Nickel: $250\times0.4096\times0.7906 =
\qty{81.0}{ppm}$. (b) $250\times0.5^{4} = \qty{15.6}{ppm}$, five times too low. (c)
The Rayleigh law assumes a constant $D$; when the crystallising assemblage changes,
$D$ jumps, and the law must be restarted from the liquid composition reached at the
change, with $F$ measured relative to that liquid.
\end{solution}

\begin{solution}{pb:b2:magmatic-differentiation:1}
\noindent\textbf{Part I.}

\textbf{1.} $\bar D_{\ce{Ni}} = 0.20\times5.0 + 0.50\times0 + 0.30\times2.0 = 1.60$.

\textbf{2.} $\bar D_{\ce{Zr}} = 0.20\times0.0007 + 0.50\times0.001 + 0.30\times0.195 =
0.00014 + 0.0005 + 0.0585 = 0.059$.

\textbf{3.} Upper zone: $\bar D_{\ce{Ni}} = 0.45\times2.0 = 0.90$; $\bar D_{\ce{Zr}} =
0.55\times0.001 + 0.45\times0.195 = 0.088$.

\textbf{4.} Nickel is compatible in the lower zone (1.60) and becomes slightly
incompatible (0.90) when olivine disappears; zirconium is incompatible in both.

\textbf{5.} Clinopyroxene: it supplies $0.0585$ of the lower-zone $0.059$, nearly
all, because it is the only mineral with a coefficient far from zero.

\textbf{6.} $(C_{\ce{Ni}}/C_{\ce{Zr}})$ changes by $F^{D_{\ce{Ni}} - D_{\ce{Zr}}}$;
per halving, $0.5^{1.60 - 0.059} = 0.5^{1.541} = 0.344$.

\noindent\textbf{Part II.}

\textbf{7.} $1.60\times150 = \qty{240}{ppm}$.

\textbf{8.} The liquid had $190/1.60 = \qty{118.75}{ppm}$; $F = (118.75/150)^{1/0.6} =
0.7917^{1.667} = 0.68$.

\textbf{9.} $80\times0.68^{-0.941} = 80\times1.438 = \qty{115}{ppm}$.

\textbf{10.} Fractional: $150\times0.68^{0.6} = 150\times0.793 = \qty{119.0}{ppm}$;
equilibrium: $150/(1.6 + 0.68\times(1 - 1.6)) = 150\times0.839 = \qty{125.8}{ppm}$.
The difference is only \qty{6}{\percent} after a third of the \omterm{def:b2:magma-emplacement:magma}{magma} has
crystallised: for a moderately \omterm{def:b2:elements-in-the-earth:partition}{compatible element} the two laws stay close early on.

\textbf{11.} $150\times0.55^{0.6} = \qty{104.8}{ppm}$.

\textbf{12.} Just below: $1.60\times104.8 = \qty{167.7}{ppm}$; just above:
$0.90\times104.8 = \qty{94.3}{ppm}$. The nickel profile of the \omterm{def:b2:magmatic-differentiation:cumulate}{cumulates} drops
abruptly by \qty{44}{\percent} at the level where olivine disappears from the
layers: a step, not a smooth curve.

\textbf{13.} Liquid: $104.8\times(0.40/0.55)^{0.90 - 1} = 104.8\times1.032 =
\qty{108.2}{ppm}$; \omterm{def:b2:magmatic-differentiation:cumulate}{cumulate} $0.90\times108.2 = \qty{97.4}{ppm}$. With $\bar D < 1$
nickel now rises slowly in the liquid and in the \omterm{def:b2:magmatic-differentiation:cumulate}{cumulates}: the profile, after its
step, climbs gently.

\noindent\textbf{Part III.}

\textbf{14.} $x_{\mathrm{ol}}(m_{\mathrm{ol},j} - d_j) + x_{\mathrm{pl}}
(m_{\mathrm{pl},j} - d_j) + x_{\mathrm{cpx}}(m_{\mathrm{cpx},j} - d_j) = 100\,(p_j -
d_j)$.

\textbf{15.} \ce{MgO}: $37.16\,x_{\mathrm{ol}} - 4.9\,x_{\mathrm{pl}} +
11.1\,x_{\mathrm{cpx}} = 410$; \ce{CaO}: $-7.3\,x_{\mathrm{ol}} +
7.06\,x_{\mathrm{pl}} + 13.2\,x_{\mathrm{cpx}} = 420$; \ce{Al2O3}:
$-14.1\,x_{\mathrm{ol}} + 17.6\,x_{\mathrm{pl}} - 11.1\,x_{\mathrm{cpx}} = 190$.

\textbf{16.} Solving (by elimination or Cramer's rule): $x_{\mathrm{ol}} = 9.03$,
$x_{\mathrm{pl}} = 30.84$, $x_{\mathrm{cpx}} = 20.32$ grams per \qty{100}{g}.

\textbf{17.} Crystals $9.03 + 30.84 + 20.32 = \qty{60.19}{g}$; liquid
\qty{39.81}{g}: $F = 0.398$.

\textbf{18.} $(4950 - 9.03\times39.19 - 30.84\times50.54 - 20.32\times51.5)/39.81 =
(4950 - 353.9 - 1558.7 - 1046.5)/39.81 = 1990.9/39.81 = \qty{50.0}{\percent}$, as in
the \omterm{def:b2:magma-emplacement:bodies}{dyke}: the oxide not used as a control is reproduced.

\textbf{19.} The crystals removed are $9.03/60.19 = \qty{15}{\percent}$ olivine,
\qty{51}{\percent} plagioclase and \qty{34}{\percent} clinopyroxene, between the
modes of the two zones. They match the lower assemblage crystallising from $F = 1$ to
0.55 ($0.45\times20 = 9.0$, $0.45\times50 = 22.5$, $0.45\times30 = 13.5$) followed by
the upper one from 0.55 to 0.40 ($0.15\times55 = 8.25$, $0.15\times45 = 6.75$): 9.0,
30.75 and 20.25 grams.

\textbf{20.} $0.30/0.75 = 0.40$.

\noindent\textbf{Part IV.}

\textbf{21.} $0.55^{0.059 - 1} = 0.55^{-0.941} = 1.755$; $80\times1.755 =
\qty{140.4}{ppm}$.

\textbf{22.} $188/140.4 = 1.339 = (F/0.55)^{0.088 - 1}$, so $F = 0.55\times1.339^{-1/
0.912} = 0.55\times0.726 = 0.399$.

\textbf{23.} $(188/80)^{1/(0.07 - 1)} = 2.35^{-1.075} = 0.399$. For \omterm{def:b2:elements-in-the-earth:partition}{incompatible
elements} $C_L/C_0 \approx 1/F$ whatever the small $D$: the answer hardly depends on
the details of $D$, which makes \omterm{def:b2:elements-in-the-earth:partition}{incompatible elements} the best gauges of $F$.

\textbf{24.} Plagioclase, which holds no iron, makes half of the crystals, and the
pyroxene and olivine together hold less iron than they hold magnesia: the liquid
grows richer in iron and poorer in magnesia without gaining silica. This iron
enrichment at nearly constant silica is the mark of the tholeiitic series of the
next chapter.

\textbf{25.} Pure cumulus crystals would hold $0.088\times188 = \qty{16.5}{ppm}$. With a
mass fraction $t$ of trapped liquid, $16.5(1 - t) + 188\,t = 47$, so $t = (47 -
16.5)/(188 - 16.5) = 30.5/171.5 = 0.18$: about \qty{18}{\percent} trapped liquid.

\textbf{26.} The three estimates, major elements (0.398), potash (0.40) and
zirconium (0.399), agree: when the uppermost layer crystallised,
\textbf{\qty[text-series-to-math=true]{40}{\percent}} of the original \omterm{def:b2:magma-emplacement:magma}{magma} was
still liquid.
\end{solution}
